\section{\texorpdfstring{Diagram-Methode für $T=0 \leftrightarrow$ Grundzustand}{Diagram-Methode für T=0 <-> Grundzustand}}
\subsection{Wechselwirkungs-Bild}

Wir zerlegen den Hamiltonoperator in
\begin{equation}
  H=H_0+H_{\mathrm{int}}.
\end{equation}
Dabei sei das Eigenwertproblem von $H_0$ lösbar.
Ziel ist es, den Einfluss von $H_{\mathrm{int}}$ systematisch zu beruecksichtigen.

\subsubsection*{Vergleich der Bilder}

\begin{center}
  \renewcommand{\arraystretch}{1.4}
  \begin{tabular}{@{}lll@{}}
    \toprule
    Bild         & Operatoren                                                                                                      & Wellenfunktion / Zustand                                                                    \\
    \midrule
    Schroedinger & $\dfrac{\partial}{\partial t}A_S=0$                                                                             & $i\hbar\,\dfrac{\partial}{\partial t}\ket{\phi_S(t)}=(H_0+H_{\mathrm{int}})\ket{\phi_S(t)}$ \\
    Heisenberg   & $A_H(t)=e^{\frac{i}{\hbar}(H_0+H_{\mathrm{int}})(t-t_0)}A_S\,e^{-\frac{i}{\hbar}(H_0+H_{\mathrm{int}})(t-t_0)}$ & $\dfrac{\partial}{\partial t}\ket{\phi_H}=0$                                                \\
    Interaction  & $A_0(t)=e^{\frac{i}{\hbar}H_0(t-t_0)}A_S\,e^{-\frac{i}{\hbar}H_0(t-t_0)}$                                       & $i\hbar\,\dfrac{\partial}{\partial t}\ket{\phi_0(t)}=H_{\mathrm{int}}^{(0)}\ket{\phi_0(t)}$ \\
    \bottomrule
  \end{tabular}
\end{center}

Zur Herleitung der Wellenfunktion im WW-Bild
benutzen wir die Picture-Unabhaengigkeit des
Erwartungswertes:
\begin{equation}
  \erw{A}
  =\bra{\phi_S(t)}A_S\ket{\phi_S(t)}
  =\underbrace{\bra{\phi_0(t)}}_{\text{Zustand im WW-Bild}}
  \,\underbrace{e^{\frac{i}{\hbar}
        H_0(t-t_0)}A_S\,e^{-\frac{i}{\hbar}H_0(t-t_0)}}_{\text{Operator }A_0(t)\text{ im WW-Bild}} \,\underbrace{\ket{\phi_0(t)}}_{\text{Zustand im WW-Bild}}.
\end{equation}
Dabei haben wir im Prinzip zwei Einsen eingefügt und benutzt, dass:
\begin{equation}
  \ket{\phi_0(t)}=e^{\frac{i}{\hbar}
      H_0(t-t_0)}\ket{\phi_S(t)}.
\end{equation}
Nun folgt die Bewegungsgleichung im WW-Bild direkt durch Ableiten:
\begin{equation}
  \begin{aligned}
    i\hbar\,\frac{\partial}{\partial t}\ket{\phi_0(t)}
     & = i\hbar\,\frac{\partial}{\partial t}\left(e^{\frac{i}{\hbar} H_0(t-t_0)}\ket{\phi_S(t)}\right)                                                  \\
     & = -H_0e^{\frac{i}{\hbar}H_0(t-t_0)}\ket{\phi_S(t)} +e^{\frac{i}{\hbar}H_0(t-t_0)}\left(i\hbar\,\frac{\partial}{\partial t}\ket{\phi_S(t)}\right) \\
     & = -H_0e^{\frac{i}{\hbar}H_0(t-t_0)}\ket{\phi_S(t)} +e^{\frac{i}{\hbar}H_0(t-t_0)}(H_0+H_{\mathrm{int}})\ket{\phi_S(t)}                           \\
     & = e^{\frac{i}{\hbar}H_0(t-t_0)}H_{\mathrm{int}}\ket{\phi_S(t)}                                                                                   \\
     & = e^{\frac{i}{\hbar}H_0(t-t_0)}H_{\mathrm{int}}e^{-\frac{i}{\hbar}H_0(t-t_0)}\ket{\phi_0(t)}                                                     \\
     & = \underbrace{e^{\frac{i}{\hbar}H_0(t-t_0)}H_{\mathrm{int}}e^{-\frac{i}{\hbar}H_0(t-t_0)}}_{\equiv\,H_{\mathrm{int}}^{(0)}(t)}\ket{\phi_0(t)}.
  \end{aligned}
\end{equation}
Genutzt wurde hier einmal die Schrödinger-Gleichung und der Fakt, dass $H_0$ mit sich selbst und mit sich selbst im Exponent $(e^{H_0})$ kommutiert.
Außerdem wurde erneut eine 1 eingefügt um aus dem Schrödinger-Zustand einen WW-Zustand zu machen.

Damit gilt asdf
\begin{equation}
  i\hbar\,\frac{\partial}{\partial t}\ket{\phi_0(t)}=H_{\mathrm{int}}^{(0)}(t)\ket{\phi_0(t)}.
\end{equation}
Hinweis: $H_0$ und $H_{\mathrm{int}}$ müssen dabei im Allgemeinen nicht kommutieren.

\subsubsection*{Formale Lösung der Bewegungsgleichung im WW-Bild}
\begin{equation}
  \label{eq:ww-eom}
  i\hbar\,\frac{\partial}{\partial t}\ket{\phi_0(t)}
  =H_{\mathrm{int}}^{(0)}(t)\ket{\phi_0(t)}.
\end{equation}
Fuer einen kleinen Zeitschritt $\delta t$:
\begin{equation}
  \begin{aligned}
    \ket{\phi_0(t+\delta t)} & = \ket{\phi_0(t)}+\delta\ket{\phi_0(t)}                                              \\
                             & =\ket{\phi_0(t)}-\frac{i}{\hbar}H_{\mathrm{int}}^{(0)}(t)\,\delta t\,\ket{\phi_0(t)} \\
                             & \approx e^{-\frac{i}{\hbar}H_{\mathrm{int}}^{(0)}(t)\,\delta t}\ket{\phi_0(t)}
    \qquad \text{Exponentialnaeerung für kleines }\delta t
  \end{aligned}
\end{equation}
Diese Allgemeine Formel können wir nun benutzen um das Intervall $[t_1,t_2]$ in vorerst endlich viele Teile Zeitintervalle $\delta t$ zu teilen.
Damit ist die gesamte Zeitentwicklung als Aneinanderreihung der $\delta t$ Zeitentwicklungen zu verstehen.
Die Zeitintervalle werden unendlich klein für eine exakte Darstellung.
\begin{equation}
  \begin{aligned}
    \ket{\phi_0(t_2)}
     & = \lim \limits_{\delta t \to 0}\cdots e^{-\frac{i}{\hbar}H_{\mathrm{int}}^{(0)}(t_{a})\delta t}e^{-\frac{i}{\hbar}H_{\mathrm{int}}^{(0)}(t_1)\delta t}\cdots\ket{\phi_0(t_1)}, \\
     & =\lim \limits_{\delta t \to 0}\mathcal T\prod_{t_i=t_1}^{t_2-\delta t}e^{-\frac{i}{\hbar}H_{\mathrm{int}}^{(0)}(t_i)\delta t}\ket{\phi_0(t_1)},
  \end{aligned}
\end{equation}
Nun ziehen wir das Produkt in die Exponentialfunktion als Summe und lassen $\delta t$ gegen 0 gehen, wodurch die Summe zu einem Integral wird.
Der Operator $\mathcal{T}$ sorgt für die richtige Zeitordnung der Exponentialfunktionen.
\begin{equation}
  \ket{\phi_0(t_2)}
  =\mathcal T\exp\!\left[-\frac{i}{\hbar}\int_{t_1}^{t_2}d\tau\,H_{\mathrm{int}}^{(0)}(\tau)\right]\ket{\phi_0(t_1)}.
\end{equation}
Schliesslich:
\begin{equation}
  \begin{aligned}
    \ket{\phi_0(t_2)}=S(t_2,t_1)\ket{\phi_0(t_1)}, \\
    S(t_2,t_1)=\mathcal T\exp\!\left[-\frac{i}{\hbar}\int_{t_1}^{t_2}d\tau\,H_{\mathrm{int}}^{(0)}(\tau)\right].
  \end{aligned}
\end{equation}

\lend{Vorl. 23.04.2026}

\subsubsection*{Eigenschaften der S-Operatoren}

\begin{equation}
  S(t_0,t)=\tilde{\mathcal T}\exp\!\left[\frac{i}{\hbar}\int_{t_0}^{t}d\tau\,H_{\mathrm{int}}^{(0)}(\tau)\right],
  \qquad t>t_0.
\end{equation}
Dabei laeuft der Operator in umgekehrter Zeitordnung, also rueckwaerts in der Zeitentwicklung. Exponent ist positiv weil die Integralgrenzen vertauscht sind.
Zum einen gilt:
\begin{equation}
  S(t_0,t)=S^{-1}(t,t_0)=S^{\dagger}(t,t_0).
\end{equation}
Damit ist $S$ unitar. Die Unitarität von $S$ koennen wir direkt zeigen:
\begin{equation}
  \begin{aligned}
    S^{-1}(t_1,t_0)S(t_1,t_0)
     & = \left[\tilde{\mathcal T}\exp\!\left(\frac{i}{\hbar}\int_{t_0}^{t_1}d\tau\,H_{\mathrm{int}}^{(0)}(\tau)\right)\right]
    \left[\mathcal T\exp\!\left(-\frac{i}{\hbar}\int_{t_0}^{t_1}d\tau\,H_{\mathrm{int}}^{(0)}(\tau)\right)\right]             \\
     & = \lim_{\delta t\to 0}
    \left[\tilde{\mathcal T}\prod_{t_i=t_0}^{t_1-\delta t}e^{\frac{i}{\hbar}H_{\mathrm{int}}^{(0)}(t_i)\delta t}\right]
    \left[\mathcal T\prod_{t_j=t_0}^{t_1-\delta t}e^{-\frac{i}{\hbar}H_{\mathrm{int}}^{(0)}(t_j)\delta t}\right]              \\
     & = 1.
  \end{aligned}
\end{equation}
Exponentielle mit Operatoren im Exponenten kürzen sich hier weg, weil die Zeitordnung genau entgegengesetzt ist.
Rechts läuft nach links, links läuft nach rechts, dadurch heben sich die Faktoren paarweise auf.

Fuer den adjungierten Operator folgt die Formel
indem man die einzelnen Operatoren adjungiert und
die Reihenfolge umkehrent.

Die naechste Eigenschaft ist die Verkettung der
Zeitentwicklung:
\begin{equation}
  S(t_3,t_2)S(t_2,t_1)=S(t_3,t_1)
  \qquad \text{fuer beliebige } t_1,t_2,t_3.
\end{equation}
Beweis fuer $t_3>t_2>t_1$:
\begin{equation}
  \begin{aligned}
    S(t_3,t_2)S(t_2,t_1)
     & = \mathcal T\exp\!\left[-\frac{i}{\hbar}\int_{t_2}^{t_3}d\tau\,H_{\mathrm{int}}^{(0)}(\tau)\right]
    \mathcal T\exp\!\left[-\frac{i}{\hbar}\int_{t_1}^{t_2}d\tau\,H_{\mathrm{int}}^{(0)}(\tau)\right]      \\
     & = \mathcal T\exp\!\left[-\frac{i}{\hbar}\int_{t_1}^{t_3}d\tau\,H_{\mathrm{int}}^{(0)}(\tau)\right] \\
     & = S(t_3,t_1).
  \end{aligned}
\end{equation}
Für die anderen Fälle (z.B. $t_2>t_3>t_1$) folgt relativ leicht:
\begin{equation}
  S(t_3,t_2)S(t_2,t_1) = S(t_3,t_1)S(t_1,t_2)S(t_2,t_1) = S(t_3,t_1).
\end{equation}

\subsubsection*{Zusammenhang zwischen Operatoren im Heisenberg- und WW-Bild}

Erwartungswert bild unabhaengig:
\begin{equation}
  \begin{aligned}
    \erw{A}
     & = \bra{\phi_0(t)}A_0(t)\ket{\phi_0(t)}                           \\
     & = \bra{\phi_0(t)}S^{\dagger}(t,t_0)A_0(t)S(t,t_0)\ket{\phi_0(t)} \\
     & = \underbrace{\bra{\phi_H}}_{\text{Heisenberg-Zustand}}
    \,\underbrace{A_H(t)}_{\text{H-Operator}}
    \,\underbrace{\ket{\phi_H}}_{\text{H-Zustand}}.
  \end{aligned}
\end{equation}
Daraus folgt unmittelbar
\begin{equation}
  A_H(t)=S^{\dagger}(t,t_0)A_0(t)S(t,t_0)=S(t_0,t)A_0(t)S(t,t_0).
\end{equation}
Für den Anfangszeitpunkt $t_0$ gilt trivialerweise
\begin{equation}
  \ket{\phi_0(t_0)}=\ket{\phi_H(t_0)},
\end{equation}
weil die Zustaende in allen Bildern am Startzeitpunkt zusammenfallen.

